\section{Discussion}
\label{sec:discussion}

The confirmatory study establishes a narrow but consequential execution result:
components of one jointly predicted action chunk can benefit from different
observation-time source ages. Because every method observes and queries at the
same controller steps, the result cannot be attributed to different feedback
frequency, query budget, or open-loop horizon. Because newest and full old20
use the same candidate cache, it also cannot be reduced to a general advantage
for fresh or delayed actions. The benefit lies in the asymmetric assignment in
this system.

This conclusion differs from the developmental question. The balanced
factorial did not support a generic incoherence penalty for mixed-source
actions. Its directional cell pattern motivated a new, fixed hypothesis, which
the untouched-state experiment then tested. Keeping that distinction visible
prevents the 62\% developmental FO result from being counted as independent
confirmation.

The result also clarifies how component-aware execution differs from related
work. DAM-VLA already demonstrates that arm and gripper can merit specialized
architectures~\cite{peng2026damvla}; our contribution is not heterogeneity
alone. TRACT routes predicted future behavior across procedural phases and
corrects arm execution~\cite{liu2026tract}; our temporal variable is instead
the past observation that generated a same-time cached prediction. TAS is the
closest source-selection formulation, but selects a complete cached action
~\cite{weng2025tas}. These distinctions make the controlled intervention, not
algorithmic complexity, the main methodological contribution.

The teacher-forced audit cautions against selecting a temporal source from
offline imitation error alone. Old predictions better match demonstrations for
both arm terms, yet the closed-loop factorial favors the fresh arm. Several
explanations are compatible with this mismatch: fresh feedback may correct
state deviation, demonstrations may reward temporal consistency differently
from the closed loop, or the separable metric may omit behaviorally important
coupling. The experiments do not distinguish these explanations. In
particular, they do not prove retained intent, non-Markovian causation, or any
general temporal-intent mechanism.

FO20 is deliberately simple. It adds no training, estimator, or policy head,
and its result suggests a future adaptive executor could predict source utility
per component. Such adaptation is not evaluated here. A learned rule would
need an independent cohort and should be compared under the same current-
observation and matched-query contract.

\subsection{Limitations}

The evidence uses one frozen ACT checkpoint and the LIBERO Object suite. It
tests one six-dimensional arm/scalar-gripper partition and confirms one source
gap, $d=20$. The study neither proves that 20 ticks is optimal nor evaluates
dynamic adaptation. It provides no causal proof that an old gripper prediction
retains intent or that non-Markovian structure produces the success gain.
There is no current confirmation on a second policy, another benchmark, a VLA,
a dexterous or bimanual action space, or a real robot. The primary cohort spans
nine tasks and paired unused states, which supports within-system credibility
but not broad generalization. Finally, success is binary; logged kinematic and
disagreement diagnostics follow treatment-dependent trajectories and are not
used to make mechanistic claims.
